\section{问题四：波动电价下的重算}

\subsection{价格信息的因果边界}

附件 4 给出的是 2025 年全年的实时电价。与负载和光伏一样，0:00 制定当天计划时无法预知当天的实际电价，因此计划必须建立在价格预测之上，而结算按实际电价进行，这正是电力市场实时结算的典型结构\cite{morales2014,huang2021}。本文沿用与前两问一致的口径，取\textbf{此前 7 天同一时段的平均电价}作为预测：
\begin{equation}
p_t^{\mathrm{fc}}=\frac1{7}\sum_{k=1}^{7}p_t^{(d-k)} .
\label{eq:pricefc}
\end{equation}

价格数据的统计特征是：附件 4 的逐日电价与附件 1 的固定日型高度相关（相关系数均值 $0.971$、中位数 $0.976$），全天均值 $0.7575$ 元/$\mathrm{kWh}$、标准差 $0.3360$ 元/$\mathrm{kWh}$，取值范围 $0.0076\sim1.7936$ 元/$\mathrm{kWh}$。式~\eqref{eq:pricefc} 的预测平均绝对误差为 $0.0827$ 元/$\mathrm{kWh}$，且几乎无系统性偏差（均值偏差 $-0.0003$）。

\subsection{模型}

\textbf{问题 4-2}（对应问题二）：0:00 用 $p^{\mathrm{fc}}$ 求解式~\eqref{eq:p2soc} 的计划并叠加安全裕度式~\eqref{eq:margin}，执行时储能严格按计划充放电，紧急购电按实际电价的 5 倍结算：
\begin{equation}
C=\sum_t p_t^{\text{real}}x_t+5\sum_t p_t^{\text{real}}E_t .
\end{equation}

\textbf{问题 4-3}（对应问题三）：光伏用式~\eqref{eq:p3expand} 的附件 3 预报，电价用式~\eqref{eq:pricefc} 的预测，按式~\eqref{eq:deviation} 做偏差结算，紧急购电同按实际电价的 5 倍。

\begin{figure}[H]
  \centering
  \includegraphics[width=0.95\textwidth]{p4_price.pdf}
  \caption{左：固定电价与实时电价的对比；右：波动电价下各方案逐月费用}
  \label{fig:p4}
\end{figure}

图~\ref{fig:p4} 左图显示实时电价在保持固定日型整体形状的同时叠加了明显的随机波动，7 日均值预测把波动抹平，因此无法捕捉当日的价格尖峰。

\subsection{结果}

\threelinetable{表 10\quad 波动电价与固定电价的结果对比（单位：万元，2025-02-01 至 2025-12-31）}{lrrr}
{项目 & 固定电价 & 波动电价 & 差额}
{对应问题二：计划购电费 & 1570.26 & 1586.77 & $+16.52$ \\
 对应问题二：紧急购电费 & 132.38 & 147.35 & $+14.97$ \\
 对应问题二：\textbf{合计} & \textbf{1702.64} & \textbf{1734.13} & $\mathbf{+31.49}$ \\
 对应问题三：购电费 & 1613.01 & 1622.37 & $+9.36$ \\
 对应问题三：偏差违约费 & 45.43 & 47.07 & $+1.64$ \\
 对应问题三：紧急购电费 & 102.12 & 113.34 & $+11.23$ \\
 对应问题三：\textbf{合计} & \textbf{1760.56} & \textbf{1782.78} & $\mathbf{+22.22}$}

\threelinetable{表 11\quad 问题四指定日期的总费用（元）}{lrr}
{日期 & 问题 4-2 & 问题 4-3}
{2025-03-20 & 46341.46 & 53060.54 \\
 2025-06-21 & 42396.35 & 41099.82 \\
 2025-09-23 & 52841.39 & 53048.42 \\
 2025-12-21 & 75219.74 & 79321.30}

\subsection{价格预测误差的代价}

为了量化价格不确定性本身的成本，本文另外计算了一条\textbf{价格完全信息}的参照线：假设 0:00 就已知当天的全部实际电价，其余设置与问题 4-2 完全相同。该方案全年费用 $17180752.62$ 元，比可实现的方案低
\begin{equation}
17341287.33-17180752.62=160534.71\ \text{元} .
\end{equation}
这就是价格预测误差在问题 4-2 口径下的全部代价，占该方案总费用的 $0.93\%$。这个数字远小于紧急购电所占的比重（$8.5\%$），说明在本数据集上\textbf{电价波动的影响小于负载与光伏预测误差的影响}：实时电价虽然波动幅度不小（标准差 $0.3360$ 元/$\mathrm{kWh}$），但它保留了固定日型的形状，而 7 日均值预测已经抓住了这个形状。

另一个佐证是：问题 4-3 的滚动修正同样有效，紧急购电由 $343800.18$ 降到 $269939.50\,\mathrm{kWh}$，降幅 $7.39$ 万 $\mathrm{kWh}$，与固定电价下的表现一致（$34.38\to26.82$ 万 $\mathrm{kWh}$）。说明滚动修正的价值来自光伏预报的改进，与价格是否波动无关。
